\section{Выводы}

В ходе выполнения лабораторной работы были изучены основные механизмы
управления процессами и межпроцессного взаимодействия в операционной
системе Linux.

Было рассмотрено создание дочернего процесса с помощью \texttt{fork()},
замена выполняемой им программы посредством семейства функций
\texttt{exec}, а также ожидание завершения дочернего процесса при помощи
\texttt{waitpid()}.

Для обмена данными между процессами был использован неименованный канал.
С помощью \texttt{dup2()} читающий конец канала был перенаправлен на
стандартный вход дочернего процесса, благодаря чему отдельная дочерняя
программа может получать передаваемые родительским процессом данные
через обычный поток \texttt{stdin}.

Также была реализована обработка ошибок системных вызовов и корректное
закрытие используемых файловых дескрипторов и файлов.

\pagebreak